\section{ЭКСПЕРИМЕНТАЛЬНОЕ СРАВНЕНИЕ NMT ПОДХОДОВ}

\subsection{Подготовка данных}

В экспериментах со всеми моделями рассматривается только русско-английская языковая пара.
Все модели обучаются переводить предложения между русским и английским языками, поэтому подготовка корпуса направлена не на построение многоязычной системы, а на получение достаточно качественных параллельных пар для одного направления перевода.
Для сравнения подходов используются три варианта набора данных разного размера, чтобы анализировать масштабируемость моделей.

Первый набор данных включает OPUS-100~\cite{opus100} и Yandex English-Russian Parallel Corpus~\cite{yandex_corpus}.
Он используется для подбора гиперпараметров, так как его размер позволяет проводить несколько серий обучения без чрезмерных вычислительных затрат.
Второй набор данных дополнительно включает ParaCrawl~\cite{paracrawl}.
Он применяется для анализа масштабируемости моделей при увеличении объема обучающих данных.
Третий набор данных строится на основе OPUS-100, Yandex Corpus, ParaCrawl и United Nations Parallel Corpus~\cite{unpc}.
Он используется для обучения и валидации лучшего из рассмотренных подходов.

Для всех трех вариантов корпуса выполняется одинаковая предварительная обработка.
На первом этапе удаляются пары с сильным расхождением длины предложений.
Пусть $l_{ru}$ и $l_{en}$ --- длины русского и английского предложений после предварительной токенизации.
Пара отбрасывается, если выполняется условие
\[
	\frac{|l_{ru} - l_{en}|}{\max(l_{ru}, l_{en})} > 0{,}3.
\]
Фильтр удаляет примеры, в которых одно предложение содержит только часть перевода другого предложения, лишний фрагмент текста или технический мусор.
При этом порог не должен быть слишком жестким, поскольку естественная длина русского и английского переводов может различаться.

Следующим этапом выполняется фильтрация по языку.
Для этого используется идея language identification, то есть автоматического определения языка текста.
В общем виде такую задачу можно описать как выбор языка с максимальной апостериорной вероятностью:
\[
	\hat{c}(s) = \arg\max_{c \in C} p(c \mid s),
\]
где $s$ --- входное предложение, $C$ --- множество возможных языков, а $c$ --- один из языковых классов.
На практике такие модели часто используют символьные признаки и статистику коротких фрагментов текста, поскольку распределения символов и символьных n-грамм хорошо различают языки даже на относительно коротких предложениях~\cite{langid}.
Пара удаляется, если для русской стороны определяется не русский язык или для английской стороны определяется не английский язык.

После этого для каждой пары вычисляется reference-free версия COMET, то есть оценка качества без эталонного перевода.
В такой постановке модель получает только два предложения пары и оценивает, насколько одно из них похоже на корректный перевод другого.
Такую постановку обычно относят к quality estimation, поскольку качество оценивается без доступа к reference-предложению~\cite{cometkiwi}.
Если обозначить русское предложение через $x$, английское предложение через $y$, а модель оценки через $f_{\phi}$, то вычисляемый score можно записать как
\[
	q = f_{\phi}(x, y).
\]
После вычисления оценки оставляются только пары, для которых выполняется условие $q > 0{,}7$.
Этот шаг удаляет часть предложений, которые прошли простые фильтры длины и языка, но все еще плохо соответствуют друг другу как переводная пара.

Дополнительно применяется аугментация с помощью LLM-модели.
Она используется не для свободного перефразирования корпуса, а для исправления локальных ошибок в уже отобранных парах.
К таким ошибкам относятся остатки разметки, технические символы, незначительные нарушения пунктуации, а также небольшие несоответствия по длине, когда одно предложение содержит начало или конец фрагмента, отсутствующий в другом предложении.
Этот этап повышает однородность корпуса без изменения основной семантики переводной пары.

После всех этапов подготовки итоговые размеры наборов данных представлены в таблице~\ref{tab:dataset-sizes}.

\begin{table}[H]
	\centering
	\caption{Размеры подготовленных наборов данных}
	\label{tab:dataset-sizes}
	\begin{tabular}{|p{0.12\textwidth}|p{0.63\textwidth}|p{0.17\textwidth}|}
		\hline
		\textbf{№} & \textbf{Состав корпуса}                                             & \textbf{Размер} \\
		\hline
		1          & OPUS-100 и Yandex Corpus                                            & 1{,}9 млн пар   \\
		\hline
		2          & OPUS-100, Yandex Corpus и ParaCrawl                                 & 4{,}7 млн пар   \\
		\hline
		3          & OPUS-100, Yandex Corpus, ParaCrawl и United Nations Parallel Corpus & 7{,}2 млн пар   \\
		\hline
	\end{tabular}
\end{table}

\subsection{Среда экспериментов и гиперпараметры}

В качестве языка реализации используется Python, а модели строятся с помощью фреймворка \texttt{torch}~\cite{pytorch}.
PyTorch удобен для экспериментальной работы с нейронными сетями: в нем можно явно описывать архитектуры, контролировать цикл обучения и использовать аппаратное ускорение для матричных операций. Кроме того, этот фреймворк широко применяется в современных исследованиях по машинному обучению.
Для корректного сравнения подходов используется единый обучающий цикл: различаются архитектуры моделей, но не базовая схема оптимизации и вычисления функции потерь.

Модели валидируются в document-level схеме перевода: весь текст параллельного корпуса подается в модель.
Качество полученного перевода оценивается по ранее обозначенным в работе метрикам: bleu, chrf и comet.

Для оптимизации используется AdamW~\cite{adamw}.
AdamW представляет собой модификацию Adam, в которой weight decay отделен от градиента основной функции потерь.
Если $g_t$ обозначает градиент на шаге $t$, то Adam оценивает экспоненциальные скользящие средние первого и второго моментов:
\[
	m_t = \beta_1 m_{t-1} + (1 - \beta_1)g_t,
	\qquad
	v_t = \beta_2 v_{t-1} + (1 - \beta_2)g_t^2.
\]
Параметры $\beta_1$ и $\beta_2$ задают скорость накопления этих средних.
После коррекции смещения параметры модели обновляются с учетом learning rate $\eta_t$ и отдельного коэффициента weight decay $\lambda$:
\[
	w_{t+1} = (1 - \eta_t\lambda)w_t
	- \eta_t\frac{\hat{m}_t}{\sqrt{\hat{v}_t} + \varepsilon}.
\]

В качестве основной функции потерь используется cross entropy loss.
Для авторегрессионной модели перевода она имеет вид:
\[
	\mathcal{L}_{CE} = -\sum_{t=1}^{n}\log p_{\theta}(y_t \mid y_{<t}, x),
\]
где $x$ --- исходное предложение, $y_t$ --- целевой токен на шаге $t$, а $y_{<t}$ --- уже известный префикс перевода.
При использовании label smoothing вместо one-hot распределения применяется сглаженное целевое распределение, что уменьшает чрезмерную уверенность модели в одном токене.
В экспериментах регуляризация включает три независимых механизма: dropout, сглаживание целевых меток и ограничение величины весов через оптимизатор.

В современных практиках обучения моделей используют не постоянный learning rate. Закон, по которому меняется learning rate, называется ``расписанием''.
Выбором расписания, как правило, устраняют эмпирически установленные недостатки различных стадий обучения модели.
Например, на самом старте обучения, когда веса модели ``далеко'' от оптимальных, градиент слишком велик, и оптимизация модели слишком нестабильна из-за проблемы ``градиентного взрыва''.
В таком случае лучше иметь меньший learning rate. Поскольку увеличение размера модели обычно усиливает нестабильность во время обучения, выставление расписания или уменьшения learning rate становится обязательным.
При этом, после нестабильного старта слишком малый learning rate иметь нежелательно, так как это замедляет процесс обучения. На поздних стадиях обучения, когда модель совершила несколько ``эпох'', происходит более детальный подбор параметров, необходимо совершать более точные шаги, а значит необходим меньший learning rate.

Одна из распространенных схем расписания learning rate решает упомянутые выше проблемы.
На первых шагах обучения применяется warmup: learning rate по линейному закону увеличивается от малого заданного значения до базового уровня.
Идея warmup широко используется в обучении Transformer-моделей~\cite{transformer}.
После завершения warmup learning rate уменьшается по косинусному закону~\cite{sgdr}, известному как ``cosine annealing'':
\[
	\eta_t = \eta_{min} + \frac{1}{2}(\eta_{max} - \eta_{min})
	\left(1 + \cos\frac{\pi(t - T_w)}{T - T_w}\right),
	\qquad t \geq T_w,
\]
где $T_w$ --- число warmup-steps, $T$ --- общее число шагов обучения, $\eta_{max}$ --- максимальное значение learning rate, а $\eta_{min}$ --- минимальное значение learning rate.
В отличие от ступенчатого уменьшения learning rate, cosine annealing не создает резких скачков в динамике оптимизации.
Описанное выше расписание требует задания двух дополнительных гиперпараметров: начальный learning rate и число warmup шагов.

Гиперпараметром можно назвать и размер словаря. Несмотря на то, что изменение размера словаря фактически меняет веса модели (поскольку меняется размер embedding слоя), сама вычислительная схема модели остается такой же.
Слишком большой словарь может излишне увеличить размер модели, что стоит учитывать при подборе этого параметра.

В каждой из рассматриваемых архитектур есть параметры, задающие размер слоев. Под скрытым слоем здесь понимается любой слой между входом нейронной сети и выходом.
Размерность входа и выхода модели, как правило, фиксированная, а вот размерность векторов между скрытыми слоями можно задавать, масштабируя слой.
Обычно между скрытыми слоями оставляют одну и ту же размерность. Сравнение моделей с разным hidden size показывает, как модель масштабируется в ширину.

Наряду с масштабированием ``в ширину'' можно изменять число одинаковых скрытых слоев, сравнивая масштабируемость модели ``в глубину''.
Изменение глубины может приводить к более существенным изменениям в процессе обучения.
Одна из причин связана с распространением градиента через последовательность слоев. Если выход слоя записать как $h_l = f_l(h_{l-1})$, то градиент по раннему представлению вычисляется по цепному правилу:
\[
	\frac{\partial \mathcal{L}}{\partial h_0}
	= \frac{\partial \mathcal{L}}{\partial h_L}
	\prod_{l=1}^{L}\frac{\partial h_l}{\partial h_{l-1}}.
\]
При увеличении числа слоев произведение якобианов становится длиннее.
Если нормы множителей в среднем меньше единицы, градиент затухает и нижние слои обновляются слабо; если больше единицы, возникает градиентный взрыв и шаги оптимизации становятся нестабильными.
Поэтому увеличение глубины влияет не только на число параметров, но и на устойчивость обучения.

Для ускорения вычислений при обучении применялся automatic mixed precision.
В качестве пониженной точности используется формат \texttt{bfloat16}, а не \texttt{float16}.
Основная причина --- сохранение 8-битной экспоненты, как у \texttt{float32}; поэтому \texttt{bfloat16} имеет близкий к \texttt{float32} диапазон представимых значений~\cite{bfloat16}.
У \texttt{float16} диапазон уже, поэтому при обучении глубоких моделей чаще возникают переполнения и исчезновение малых градиентов, а также может потребоваться дополнительное масштабирование loss функции.
Формат \texttt{bfloat16} имеет меньшую точность мантиссы, но для обучения нейронных сетей это обычно менее критично, чем устойчивый диапазон значений.
Поэтому \texttt{bfloat16} уменьшает потребление памяти и ускоряет матричные операции без существенного изменения гиперпараметров обучения.

\subsection{Схема подбора гиперпараметров}

Полный перебор гиперпараметров для всех рассматриваемых архитектур практически невозможен. Даже при небольшом числе значений для learning rate, размера batch, hidden size, числа слоев, параметров регуляризации и расписания обучения количество запусков быстро становится слишком большим. При этом каждый запуск требует обучения модели на корпусе значительного размера, поэтому схема экспериментов должна учитывать ограничение вычислительных ресурсов.

В работе используется допущение, что модели разных архитектур имеют часть общих свойств обучения. Например, слишком большой learning rate может приводить к нестабильной оптимизации как для Transformer, так и для Mamba, а чрезмерный dropout может ухудшать сходимость независимо от конкретного блока. Поэтому глобальные параметры, подобранные на одной модели, рассматриваются как приближенно подходящие и для других моделей. Такое допущение не означает, что найденные значения являются строго оптимальными для каждой архитектуры, но позволяет сократить число экспериментов до практически выполнимого набора.

Второе допущение состоит в том, что большинство гиперпараметров можно подбирать последовательно. В этом случае сначала фиксируется базовая конфигурация, затем изменяется один параметр или небольшая группа связанных параметров, после чего выбранное значение используется в следующих запусках. Такой подход не гарантирует нахождение глобального оптимума, однако позволяет оценить влияние отдельных решений без полного перебора всех комбинаций.

Некоторые параметры нельзя считать независимыми. Learning rate связан с размером batch: при изменении числа примеров в одном шаге меняется шум градиента и масштаб обновлений. Параметры warmup также рассматриваются совместно, поскольку число warmup-steps и начальное значение learning rate определяют одну и ту же начальную фазу оптимизации. Hidden size, число слоев и размер обучающего набора связаны через емкость модели: увеличение глубины или ширины без достаточного объема данных может усиливать переобучение, а слишком малая модель может не использовать преимущества большего корпуса. Поэтому такие параметры подбираются не по отдельности, а как связанные группы.

Значения loss корректно сравнивать только для запусков с одинаковыми конфигурациями словарей и одинаковым значением label smoothing. Размер и состав словаря напрямую определяют пространство классов в cross entropy loss, а label smoothing изменяет целевое распределение вероятностей. Поэтому при различии этих параметров loss отражает не только качество модели, но и изменение самой функции потерь. На практике увеличение label smoothing искусственно завышает значение loss, даже если качество перевода по внешним метрикам не ухудшается.

\subsection{Эксперименты с токенизатором}

При обучении токенизаторов использовалось предварительное разбиение, в котором цифры и знаки препинания отделялись от соседних буквенных фрагментов, так как числа и пунктуация часто должны переноситься в перевод почти без изменения. Это было подтверждено экспериментально: добавление такой предобработки улучшило точность переноса чисел и знаков препинания в переведенный текст.

Размер словаря выбирался по моменту, после которого увеличение числа подслов почти не отражалось на внешних метриках. Для небольшого корпуса словарь больше 36 тыс. токенов оказывался избыточным: новые элементы встречались редко, но сразу увеличивали embedding-слой и выходную проекцию. На среднем наборе данных рабочим оказался уровень около 48 тыс. токенов. Для большого корпуса граница смещалась примерно к 64 тыс.; редкие подслова в этом случае получают больше наблюдений, и разрежение статистики становится менее заметным.

Для Transformer encoder-decoder схемы был сделан дополнительный запуск с BPE. Обучались две двуязычные модели со словарем 48 тыс. токенов. Конфигурация базовой encoder-decoder модели приведена в приложении~А. Обучение шло одну эпоху на датасете среднего размера; архитектура, optimizer и расписание learning rate не менялись. В таблицу не включен loss, поскольку при другом словаре меняется число классов и распределение частот токенов.

\begin{table}[H]
	\centering
	\caption{Сравнение unigram и BPE для базовой Transformer encoder-decoder модели}
	\label{tab:tokenizer-comparison}
	\begin{tabular}{|p{0.32\textwidth}|c|c|c|}
		\hline
		\textbf{Токенизатор} & \textbf{BLEU} & \textbf{chrF} & \textbf{COMET} \\
		\hline
		unigram              & 20{,}78       & 51{,}04       & 0{,}7385       \\
		\hline
		BPE                  & 21{,}07       & 51{,}47       & 0{,}7383       \\
		\hline
	\end{tabular}
\end{table}

BPE дал +0{,}29 BLEU и +0{,}43 chrF при практически совпадающем COMET.
Такой результат не распространился на все архитектуры, но для Transformer encoder-decoder BPE оказался лучшим вариантом.
В межархитектурных сравнениях далее будет использоваться unigram, так как он дает более стабильные результаты в среднем по альтернативным (LSTM, Mamba) архитектурам.

\subsection{Эксперименты с параметрами обучения}

В следующих экспериментах использовалась модель, уже заданная в экспериментах с токенизатором (``небольшой'' transformer encoder-decoder).
Эксперименты также проводились на протяжении одной эпохи датасета средней конфигурации.

\paragraph{Расписание learning rate.}
На основе нескольких первых экспериментов, до полноценного подбора гиперпараметров, был выбран единый формат расписания learning rate - linear warmup + cosine annealing с фиксированными гиперпараметрами.
В warmup фазе было решено использовать learning rate равный 1\% от базового значения, на протяжении одной пятой эпохи.
Минимальное значение при annealing выставлено равным 10\% от базового learning rate.
Примерный вид использованного расписания показан на рисунке~\ref{fig:lr-schedule}, фрагмент создания scheduler в PyTorch приведен в приложении~А.

\begin{figure}[H]
	\centering
	\begin{tikzpicture}
		\begin{axis}[
				width=0.75\textwidth,
				height=0.38\textwidth,
				xlabel={Доля обучения},
				ylabel={Отношение к базовому learning rate},
				xmin=0, xmax=1,
				ymin=0, ymax=1.05,
				grid=major,
				xtick={0,0.2,1},
				xticklabels={$0$, warmup, $1$},
				ytick={0.01,0.1,1},
				yticklabels={$0{,}01$, $0{,}1$, $1$},
			]
			\addplot[thick, domain=0:0.2, samples=2] {0.01 + (1 - 0.01) * x / 0.2};
			\addplot[thick, domain=0.2:1, samples=100] {0.1 + 0.5 * (1 - 0.1) * (1 + cos(deg((x - 0.2) / 0.8 * pi)))};
		\end{axis}
	\end{tikzpicture}
	\caption{Схематичный вид расписания learning rate}
	\label{fig:lr-schedule}
\end{figure}

После фиксации формы расписания сравнивались несколько базовых значений learning rate.
Основное сравнение выполнялось при достаточно большом batch size, равном 100.
Слишком малое значение $3 \cdot 10^{-5}$ приводило к медленному обучению: за одну эпоху модель оставалась на низком уровне качества.
Значения $3 \cdot 10^{-4}$ и $3 \cdot 10^{-3}$ уже давали рабочую динамику, при этом в данном коротком запуске более высокий learning rate оказался лучше по всем внешним метрикам.
На более длинных запусках была замечена закономерность при использовании слишком больших learning rate: функция потерь, несмотря на быстрое уменьшение в начале, также быстрее выходила на плато, из-за чего конечный результат был хуже.

\begin{table}[H]
	\centering
	\caption{Влияние базового learning rate на качество перевода}
	\label{tab:lr-comparison}
	\begin{tabular}{|p{0.24\textwidth}|c|c|c|c|}
		\hline
		\textbf{Learning rate} & \textbf{BLEU} & \textbf{chrF} & \textbf{COMET} & \textbf{Loss} \\
		\hline
		$3 \cdot 10^{-5}$      & 6{,}35        & 30{,}49       & 0{,}4566       & 4{,}312       \\
		\hline
		$3 \cdot 10^{-4}$      & 20{,}96       & 51{,}26       & 0{,}7384       & 3{,}066       \\
		\hline
		$3 \cdot 10^{-3}$      & 22{,}97       & 52{,}93       & 0{,}7658       & 2{,}919       \\
		\hline
	\end{tabular}
\end{table}

При уменьшении batch size соотношение изменилось: значение $3 \cdot 10^{-4}$ оказалось устойчивее, чем $3 \cdot 10^{-3}$. Дополнительное сравнение приведено в таблице~\ref{tab:lr-small-batch-comparison}.

\begin{table}[H]
	\centering
	\caption{Влияние learning rate при меньшем batch size}
	\label{tab:lr-small-batch-comparison}
	\begin{tabular}{|p{0.24\textwidth}|c|c|c|c|}
		\hline
		\textbf{Learning rate} & \textbf{BLEU} & \textbf{chrF} & \textbf{COMET} & \textbf{Loss} \\
		\hline
		$3 \cdot 10^{-5}$      & 15{,}84       & 45{,}72       & 0{,}6731       & 3{,}487       \\
		\hline
		$3 \cdot 10^{-4}$      & 21{,}64       & 52{,}04       & 0{,}7506       & 3{,}018       \\
		\hline
		$3 \cdot 10^{-3}$      & 20{,}88       & 50{,}92       & 0{,}7337       & 3{,}141       \\
		\hline
	\end{tabular}
\end{table}

Такое поведение можно объяснить через стохастическую природу mini-batch обучения. Градиент, вычисленный по batch, является шумной оценкой полного градиента, причем дисперсия этой оценки уменьшается при росте batch size. Если batch size уменьшается, то при том же learning rate возрастает относительная величина случайной составляющей шага оптимизации. В работах по связи learning rate и batch size этот эффект описывается через noise scale, который для малых batch относительно размера датасета приблизительно пропорционален отношению learning rate к batch size~\cite{batch_size_lr}. Поэтому уменьшение batch size может требовать уменьшения learning rate: в противном случае шаги становятся более шумными и модель быстрее попадает в менее устойчивую область оптимизации.

\paragraph{Параметры AdamW.}
Для AdamW отдельно проверялись коэффициенты экспоненциального сглаживания первого и второго моментов градиента. При изменении $\beta_1$ значение $\beta_2$ оставалось фиксированным, а при изменении $\beta_2$ фиксировалось $\beta_1$.

\begin{table}[H]
	\centering
	\caption{Влияние $\beta_1$ в AdamW на качество перевода}
	\label{tab:beta1-comparison}
	\begin{tabular}{|p{0.24\textwidth}|c|c|c|c|}
		\hline
		\textbf{$\beta_1$} & \textbf{BLEU} & \textbf{chrF} & \textbf{COMET} & \textbf{Loss} \\
		\hline
		0{,}85             & 21{,}48       & 51{,}61       & 0{,}7409       & 3{,}063       \\
		\hline
		0{,}90             & 20{,}82       & 50{,}87       & 0{,}7363       & 3{,}076       \\
		\hline
		0{,}95             & 21{,}06       & 51{,}29       & 0{,}7396       & 3{,}057       \\
		\hline
	\end{tabular}
\end{table}

Различия между значениями $\beta_1$ получились небольшими. Наилучший BLEU, chrF и COMET в этой группе дал вариант $0{,}85$, но преимущество над $0{,}95$ не выглядит большим.
Поэтому $\beta_1$ можно считать, вероятно, менее чувствительным параметром, чем learning rate.

\begin{table}[H]
	\centering
	\caption{Влияние $\beta_2$ в AdamW на качество перевода}
	\label{tab:beta2-comparison}
	\begin{tabular}{|p{0.24\textwidth}|c|c|c|c|}
		\hline
		\textbf{$\beta_2$} & \textbf{BLEU} & \textbf{chrF} & \textbf{COMET} & \textbf{Loss} \\
		\hline
		0{,}95             & 18{,}05       & 47{,}72       & 0{,}6820       & 3{,}342       \\
		\hline
		0{,}98             & 19{,}97       & 49{,}84       & 0{,}7166       & 3{,}180       \\
		\hline
		0{,}999            & 21{,}20       & 51{,}27       & 0{,}7383       & 3{,}064       \\
		\hline
	\end{tabular}
\end{table}

Для $\beta_2$ зависимость выражена сильнее: увеличение коэффициента до $0{,}999$ улучшило все внешние метрики и одновременно снизило loss.
Вероятная причина состоит в более стабильной оценке второго момента градиента.
Результат интересен в первую очередь тем, что большинство современных больших языковых моделей обучается с $\beta_2 = 0{,}95$, но это значение не является оптимальным в проведенных экспериментах.

\paragraph{Weight decay.}
Weight decay применялся не ко всем параметрам модели: из него исключались bias-параметры, embedding-слои и normalization-слои~\cite{l2_normalization_decay}.
Embedding матрицу можно рассматривать как таблицу значений, сопоставляющую одному токену его векторное представление.
Если в процессе градиентного спуска применить decay к этой таблице, векторы, соответствующие редким токенам, будут задействованы редко, и градиент по ним ``протекать'' не будет. Таким образом, decay будет просто ``тянуть'' векторы редких токенов к нулю, и только иногда они будут обновляться:
\[
	e_{i,t+1} =
	\begin{cases}
		(1 - \eta_t \lambda)e_{i,t} - \eta_t u_{i,t}, & i \in B_t,    \\
		(1 - \eta_t \lambda)e_{i,t},                  & i \notin B_t,
	\end{cases}
\]
где $e_i$ -- embedding-вектор токена $i$, $B_t$ -- множество токенов в текущем mini-batch, $\eta_t$ -- learning rate, $\lambda$ -- weight decay, а $u_{i,t}$ -- градиентное обновление для этого embedding-вектора.
Если токен $i$ редкий и не встречается в течение $K$ шагов подряд, то для него $i \notin B_t$ на каждом из этих шагов. В таком случае embedding-вектор не получает содержательного градиентного обновления и изменяется только за счет decay:
\[
	e_{i,t+K}
	= e_{i,t}\prod_{s=t}^{t+K-1}(1 - \eta_s\lambda)
	\approx e_{i,t}\exp\left(-\lambda\sum_{s=t}^{t+K-1}\eta_s\right).
\]
При положительных $\eta_s$ и $\lambda$ множитель перед $e_{i,t}$ становится меньше единицы, поэтому норма embedding-вектора постепенно уменьшается. Когда редкий токен снова появляется в batch, его вектор уже может быть ближе к нулю, и одно градиентное обновление потенциально будет снова обучать модель тому, что она уже знала.
Для слоев нормализации weight decay форсирует модель работать с низкочастотными сигналами, так как значение scale параметра нормализации стягивается к нулю. Аналогичное объяснение можно применить и к bias параметрам линейных слоев -- decay тянет значение сдвига к нулю, убирая всю пользу bias параметров в целом.
Фрагмент реализации такого разделения параметров на уровне оптимизатора приведен в приложении~А.

\begin{table}[H]
	\centering
	\caption{Влияние weight decay на качество перевода}
	\label{tab:weight-decay-comparison}
	\begin{tabular}{|p{0.24\textwidth}|c|c|c|c|}
		\hline
		\textbf{Weight decay} & \textbf{BLEU} & \textbf{chrF} & \textbf{COMET} & \textbf{Loss} \\
		\hline
		0{,}005               & 20{,}82       & 50{,}88       & 0{,}7364       & 3{,}078       \\
		\hline
		0{,}01                & 20{,}85       & 50{,}94       & 0{,}7371       & 3{,}071       \\
		\hline
		0{,}1                 & 21{,}03       & 51{,}18       & 0{,}7390       & 3{,}055       \\
		\hline
	\end{tabular}
\end{table}

Различия между значениями weight decay получились небольшими. Наилучшие значения внешних метрик дал вариант $0{,}1$, однако отрыв от $0{,}01$ и $0{,}005$ остается малым. Поэтому этот параметр в проверенном диапазоне выглядит менее чувствительным, чем learning rate или $\beta_2$.

\paragraph{Dropout.}
Dropout проверялся как основной параметр регуляризации модели.
В коротком запуске на датасете среднего размера более сильное зануление активаций не дало преимущества: качество постепенно снижалось при переходе от $0{,}05$ к $0{,}2$.

\begin{table}[H]
	\centering
	\caption{Влияние dropout на качество перевода}
	\label{tab:dropout-comparison}
	\begin{tabular}{|p{0.24\textwidth}|c|c|c|c|}
		\hline
		\textbf{Dropout} & \textbf{BLEU} & \textbf{chrF} & \textbf{COMET} & \textbf{Loss} \\
		\hline
		0{,}05           & 21{,}40       & 51{,}21       & 0{,}7380       & 3{,}058       \\
		\hline
		0{,}10           & 20{,}69       & 50{,}91       & 0{,}7348       & 3{,}071       \\
		\hline
		0{,}20           & 20{,}10       & 50{,}51       & 0{,}7311       & 3{,}093       \\
		\hline
	\end{tabular}
\end{table}

Однако модели большего размера ($8$ слоев encoder и decoder) без увеличения размера датасета лучше показывали себя на конфигурации с $dropout = 0{,}1$.

\begin{table}[H]
	\centering
	\caption{Влияние dropout для модели большего размера}
	\label{tab:dropout-large-model-comparison}
	\begin{tabular}{|p{0.24\textwidth}|c|c|c|c|}
		\hline
		\textbf{Dropout} & \textbf{BLEU} & \textbf{chrF} & \textbf{COMET} & \textbf{Loss} \\
		\hline
		0{,}05           & 23{,}74       & 53{,}31       & 0{,}7724       & 2{,}842       \\
		\hline
		0{,}10           & 24{,}18       & 53{,}86       & 0{,}7791       & 2{,}811       \\
		\hline
		0{,}20           & 22{,}96       & 52{,}74       & 0{,}7638       & 2{,}889       \\
		\hline
	\end{tabular}
\end{table}

Эксперимент подтверждает эмпирическое суждение о том, что модели большего размера сильнее подвержены переобучению. Этот факт учитывался при проведении финальных экспериментов.

\paragraph{Label smoothing.}
Label smoothing проверялся на двух одинаковых Transformer encoder-decoder конфигурациях; менялось только значение коэффициента сглаживания. Конфигурация модели приведена в приложении~А. В первом запуске использовалось $0{,}05$, во втором --- $0{,}1$.

При сглаживании часть вероятностной массы переносится с правильного токена на остальные токены словаря.
Loss между такими запусками напрямую не сравнивался: при изменении label smoothing меняется целевое распределение, а значит и сама численная шкала функции потерь. Поэтому ниже приведены только внешние метрики качества перевода.

\begin{table}[H]
	\centering
	\caption{Влияние label smoothing на качество перевода}
	\label{tab:label-smoothing-comparison}
	\begin{tabular}{|p{0.32\textwidth}|c|c|c|}
		\hline
		\textbf{Label smoothing} & \textbf{BLEU} & \textbf{chrF} & \textbf{COMET} \\
		\hline
		$0{,}05$                 & 31{,}68       & 60{,}19       & 0{,}8474       \\
		\hline
		$0{,}1$                  & 32{,}33       & 60{,}61       & 0{,}8497       \\
		\hline
	\end{tabular}
\end{table}

При $0{,}1$ все три внешние метрики оказались выше: BLEU на 0{,}65, chrF на 0{,}42, COMET на 0{,}0023. Разница невелика, но направлена одинаково по всем метрикам, поэтому $0{,}1$ был взят как рабочее значение для последующих запусков.

\paragraph{GRPO-дообучение.}
Для encoder-decoder модели был выполнен короткий этап GRPO-дообучения. За основу бралась уже обученная Transformer encoder-decoder конфигурация; архитектура при этом не менялась. Конфигурация модели и параметры GRPO-дообучения приведены в приложении~А. Reward задавался как взвешенная сумма нормированных автоматических метрик качества перевода:
\[
	R = \frac{1}{3} \cdot \frac{\operatorname{chrF}}{60} +
	\frac{2}{3} \cdot \frac{\operatorname{COMET}}{0{,}845}.
\]
COMET получил больший вес, поскольку reward должен был учитывать не только поверхностные совпадения с reference, но и смысловую близость перевода. chrF оставлен как более простой и устойчивый символьный сигнал. BLEU в reward не включался: для коротких предложений он часто дает шумную оценку и может поощрять слишком буквальное совпадение $n$-грамм.

Для каждого исходного предложения генерировалась группа из четырех вариантов перевода. Из-за дополнительной генерации (семплирования) GRPO-шаг значительно дороже обычного шага supervised-обучения: модель не только считает loss по готовому target, но и сначала порождает несколько кандидатов.
Поэтому за время дообучения было обработано только около 50 тыс. обучающих примеров.

Сравнение базовой модели и модели после GRPO-дообучения приведено в таблице~\ref{tab:grpo-comparison}.

\begin{table}[H]
	\centering
	\caption{Сравнение базовой модели и GRPO-дообучения}
	\label{tab:grpo-comparison}
	\begin{tabular}{|p{0.42\textwidth}|c|c|c|}
		\hline
		\textbf{Модель}       & \textbf{BLEU} & \textbf{chrF} & \textbf{COMET} \\
		\hline
		Базовая модель        & 32{,}12       & 60{,}37       & 0{,}8497       \\
		\hline
		После GRPO-дообучения & 32{,}12       & 60{,}71       & 0{,}8527       \\
		\hline
	\end{tabular}
\end{table}

После GRPO-дообучения BLEU остался на том же уровне, chrF вырос на 0{,}34, COMET -- на 0{,}0030. Улучшение произошло только в тех метриках, что явно входили в формулу reward, поэтому результат скорее похож на подстройку модели под выбранную награду, чем на общее улучшение перевода.
При этом затраченное время было сопоставимо с заметной частью supervised-обучения.
Таким образом более выгодным выглядит улучшение данных и supervised этапа; для GRPO, вероятно, нужна более экономичная схема генерации или другая reward-функция.
Ожидания от этого этаба обучения были выше и не оправдались.

\subsection{Эксперименты с архитектурами моделей}

Архитектуры сравнивались на схожем размере по числу параметров. Точного равенства добиться сложно: у encoder-decoder и decoder-only моделей по-разному распределяются параметры между embedding, блоками последовательности и выходной проекцией. Тем не менее большинство запусков попадало в диапазон примерно от 90 до 130 млн обучаемых параметров. Для обучения использовался датасет среднего размера с подобранными ранее параметрами обучения.
Результаты обучения на основе 10 полных эпох представлены в таблице~\ref{tab:architecture-comparison}.

\begin{table}[H]
	\centering
	\caption{Сравнение архитектур моделей NMT}
	\label{tab:architecture-comparison}
	\begin{tabularx}{\textwidth}{|p{0.29\textwidth}|X|c|c|c|}
		\hline
		\textbf{Архитектура}         & \textbf{Схема}                                    & \textbf{BLEU} & \textbf{chrF} & \textbf{COMET} \\
		\hline
		LSTM Encoder-Decoder         & encoder-decoder, LSTM с attention                 & 26{,}18       & 55{,}9        & 0{,}813        \\
		\hline
		Transformer Encoder-Decoder  & encoder-decoder, self-attention и cross-attention & 29{,}34       & 58{,}8        & 0{,}836        \\
		\hline
		Mamba Encoder-Decoder гибрид & encoder-decoder, Mamba-блоки и cross-attention    & 28{,}94       & 58{,}6        & 0{,}833        \\
		\hline
		Transformer Decoder          & decoder-only, causal self-attention               & 27{,}87       & 57{,}3        & 0{,}824        \\
		\hline
		Mamba Decoder                & decoder-only, каузальные Mamba-блоки              & 18{,}66       & 49{,}4        & 0{,}706        \\
		\hline
	\end{tabularx}
\end{table}

LSTM Encoder-Decoder уступил Transformer-вариантам по всем трем метрикам; разрыв по BLEU превышает 3 пункта. Результат в целом ожидаем -- Transformer архитектура со временем показала свое преимущество над рекуррентными подходами.
Однако, результат LSTM модели все еще удовлетворительный, модель смогла выдавать адекватный и различимый перевод на текстах малого размера.

Среди encoder-decoder запусков Transformer и Mamba-гибрид оказались близки. Transformer выше на 0{,}40 BLEU, 0{,}2 chrF и 0{,}003 COMET, то есть разрыв небольшой относительно общей разницы между архитектурными группами. Такой результат разумнее трактовать осторожно: Mamba-блоки в этой конфигурации не разрушили качество при замене части self-attention, но и не дали заметного преимущества. Возможно, решающим для обеих моделей здесь был не столько сам блок внутри encoder и decoder, сколько сохраненный cross-attention между исходной и целевой последовательностями.

Decoder-only модели на таком масштабе размеров оказались хуже encoder-decoder вариантов.
Поскольку Decoder модели довольно просты по своему устройству и лучше всего подходят для общей задачи языкового моделирования, их можно обучить на большом корпусе любых текстов в целом, в совсем другом масштабе параметров, чем исследуемые в работе.
Такая модель, содержащая общие знания о языке, вероятно сможет ценой большей вычислительно нагрузки показать результат лучший, чем encoder-decoder модель.
Именно по этой причине современные LLM все чаще используются для задачи NMT -- их проще масштабировать и дообучать под конкретный домен перевода.
Однако в условиях ограниченных вычислительных ресурсов можно заключить, что encoder-decoder модели показывают себя лучше.

Mamba Decoder получил самый низкий результат.
Здесь одновременно отсутствуют отдельный encoder, cross-attention и попарный self-attention между всеми позициями.
Поэтому модель должна удерживать сведения об исходном предложении во внутреннем состоянии каузального блока.
Для языкового продолжения текста такой режим может быть естественным, но для перевода он предъявляет более жесткие требования к сохранению деталей исходной последовательности, что на таком размере модели скорее всего, невозможно.

В качестве оптимальной была выбрана Transformer Encoder-Decoder архитектура. В сравнении encoder-decoder Transformer дал лучший набор метрик и оставил понятные параметры для последующей настройки --- позиционное представление, размер словаря, регуляризацию, глубину и hidden size. Ключевой фрагмент реализации модели приведен в приложении~А.

\subsection{Эксперименты с параметрами моделей}

\paragraph{Позиционное представление.}
Для оценки влияния способа задания позиции сравнивались две Transformer encoder-decoder модели сопоставимого размера. В первой модели использовалось absolute positional encoding, во второй --- rotary positional embeddings. Остальные основные архитектурные параметры совпадали: число слоев encoder и decoder, hidden size, число attention-голов, размер feed-forward блока и dropout. Конфигурации моделей приведены в приложении~А. Для обучения был использован датасет большого размера.

Сравнение выполнялось по лучшим значениям метрик на validation-наборе, было совершено 10 полных эпох обучения. Результаты представлены в таблице~\ref{tab:pe-rope-comparison}.

\begin{table}[H]
	\centering
	\caption{Сравнение positional encoding и RoPE}
	\label{tab:pe-rope-comparison}
	\begin{tabular}{|p{0.32\textwidth}|c|c|c|c|}
		\hline
		\textbf{Позиционное представление} & \textbf{BLEU} & \textbf{chrF} & \textbf{COMET} & \textbf{Loss} \\
		\hline
		Absolute PE                        & 31{,}22       & 59{,}78       & 0{,}8450       & 1{,}931       \\
		\hline
		RoPE                               & 31{,}59       & 59{,}84       & 0{,}8484       & 1{,}914       \\
		\hline
	\end{tabular}
\end{table}

Разница в пользу RoPE небольшая: сильнее всего изменились BLEU и COMET, chrF почти не сдвинулся. Тем не менее направление изменения одинаковое для всех метрик, поэтому RoPE был оставлен в следующих Transformer encoder-decoder запусках. Возможное объяснение связано с тем, что позиция в RoPE входит не как добавка к embedding, а непосредственно в скалярные произведения query и key, от которых зависят attention-веса.

\paragraph{Направление перевода.}
Также сравнивались две encoder-decoder модели в однонаправленной и двуязычной постановках. В однонаправленной постановке модель обучалась переводить только из русского языка в английский. В двуязычной постановке одна модель обучалась обоим направлениям перевода с указанием целевого языка во входной последовательности. Для сравнения использовались модели с одинаковыми основными архитектурными параметрами: Transformer encoder-decoder с RoPE, 8 слоев encoder, 8 слоев decoder, hidden size 768, 12 attention-голов, feed-forward multiplier 3 и dropout 0{,}05.

Результаты сравнения приведены в таблице~\ref{tab:unidirectional-bilingual-comparison}.

\begin{table}[H]
	\centering
	\caption{Сравнение однонаправленной и двуязычной encoder-decoder постановки}
	\label{tab:unidirectional-bilingual-comparison}
	\begin{tabular}{|p{0.32\textwidth}|c|c|c|c|}
		\hline
		\textbf{Постановка обучения} & \textbf{BLEU} & \textbf{chrF} & \textbf{COMET} & \textbf{Loss} \\
		\hline
		Однонаправленная             & 31{,}59       & 59{,}84       & 0{,}8484       & 1{,}914       \\
		\hline
		Двуязычная                   & 32{,}05       & 60{,}35       & 0{,}8493       & 1{,}819       \\
		\hline
	\end{tabular}
\end{table}

Двуязычная постановка в этом сравнении оказалась выше по BLEU, chrF и COMET, а validation loss также снизился. Возможная причина --- более полное использование параллельного корпуса: обе стороны выступают и как источник, и как целевой текст. Параметры encoder и decoder при этом остаются общими, так что модель получает больше разнообразных обучающих примеров без изменения архитектуры.

\paragraph{Масштабирование размера.}
В рамках поставленных ранее экспериментов была получена ожидаемая связь размера модели и размера обучающего набора данных. Увеличение числа слоев и hidden size улучшало качество только до тех пор, пока обучающий корпус давал достаточно примеров для настройки дополнительных параметров. После этого более крупная модель начинала обучаться медленнее, а прирост внешних метрик становился небольшим или исчезал.

Ограничение также зависело от архитектуры. LSTM encoder-decoder модель было трудно масштабировать выше 8 слоев encoder и decoder: обучение становилось заметно медленнее, а в отдельных запусках появлялась нестабильность, похожая на градиентный взрыв.

Для Transformer encoder-decoder и гибридной Mamba encoder-decoder архитектур масштабирование выглядело более устойчивым, удалось установить оптимальные варианты конфигурации моделей в зависимости от объема обучающих данных, приведенные в таблице~\ref{tab:model-size-data-scale}.

\begin{table}[H]
	\centering
	\caption{Практический предел размера Transformer/Mamba encoder-decoder модели}
	\label{tab:model-size-data-scale}
	\begin{tabular}{|p{0.32\textwidth}|c|c|}
		\hline
		\textbf{Размер датасета} & \textbf{Число слоев encoder/decoder} & \textbf{Hidden size} \\
		\hline
		Малый                    & 6/6                                  & 512                  \\
		\hline
		Средний                  & 8/8                                  & 768                  \\
		\hline
		Большой                  & 12/12                                & 1024                 \\
		\hline
	\end{tabular}
\end{table}

При дальнейшем увеличении размера модель получала слишком большую емкость относительно объема данных, поэтому качество на validation-наборе росло слабо. Для последующего итогового обучения был выбран промежуточный вариант Transformer encoder-decoder: 8 слоев encoder, 8 слоев decoder и hidden size 768. Он соответствовал датасету среднего и большого размера, но не требовал столь больших вычислительных затрат, как конфигурация 12/12 с hidden size 1024.

\subsection{Обучение лучшей конфигурации}

После подбора гиперпараметров была обучена итоговая двуязычная модель перевода. В качестве архитектуры используется Transformer encoder-decoder с rotary positional embeddings. Модель содержит около 230 млн параметров и обучается на наибольшем подготовленном наборе данных в течение 10 эпох.
Основные параметры итогового обучения приведены в таблице~\ref{tab:final-model-hparams}.

\begin{table}[H]
	\centering
	\caption{Гиперпараметры итоговой модели}
	\label{tab:final-model-hparams}
	\begin{tabularx}{\textwidth}{|p{0.46\textwidth}|X|}
		\hline
		\textbf{Параметр}              & \textbf{Значение}                      \\
		\hline
		Архитектура                    & Transformer encoder-decoder с RoPE     \\
		\hline
		Число параметров               & 230 млн                                \\
		\hline
		Языковая постановка            & двуязычная модель ru--en               \\
		\hline
		Число слоев encoder и decoder  & 8 и 8                                  \\
		\hline
		Hidden size                    & 768                                    \\
		\hline
		Число attention-голов          & 12                                     \\
		\hline
		Число эпох                     & 10                                     \\
		\hline
		Batch size                     & 48                                     \\
		\hline
		Learning rate                  & 0{,}0003                               \\
		\hline
		Минимальный learning rate      & 0{,}00001                              \\
		\hline
		Начальный learning rate warmup & $0{,}01 \cdot lr$                      \\
		\hline
		Параметры AdamW                & $\beta_1 = 0{,}9$, $\beta_2 = 0{,}999$ \\
		\hline
		Label smoothing                & 0{,}1                                  \\
		\hline
		Dropout                        & 0{,}1                                  \\
		\hline
		Weight decay                   & 0{,}05                                 \\
		\hline
	\end{tabularx}
\end{table}

Во время обучения после каждой десятой части эпохи вычислялись значения validation loss, BLEU, chrF и COMET. Динамика BLEU, chrF и COMET на validation-наборе показана на рисунках~\ref{fig:final-validation-bleu}--\ref{fig:final-validation-comet}.

\begin{figure}[H]
	\centering
	\begin{tikzpicture}
		\begin{axis}[
				width=0.9\textwidth,
				height=0.42\textwidth,
				xlabel={Шаг валидации},
				ylabel={BLEU},
				xmin=1,xmax=100,
				grid=both,
				major grid style={gray!30},
				minor grid style={gray!15},
			]
			\addplot[blue, thick] table[x=step,y=bleu,col sep=comma] {csv/final_validation.csv};
		\end{axis}
	\end{tikzpicture}
	\caption{Динамика BLEU итоговой модели на validation-наборе}
	\label{fig:final-validation-bleu}
\end{figure}

\begin{figure}[H]
	\centering
	\begin{tikzpicture}
		\begin{axis}[
				width=0.9\textwidth,
				height=0.42\textwidth,
				xlabel={Шаг валидации},
				ylabel={chrF},
				xmin=1,xmax=100,
				grid=both,
				major grid style={gray!30},
				minor grid style={gray!15},
			]
			\addplot[red, thick] table[x=step,y=chrf,col sep=comma] {csv/final_validation.csv};
		\end{axis}
	\end{tikzpicture}
	\caption{Динамика chrF итоговой модели на validation-наборе}
	\label{fig:final-validation-chrf}
\end{figure}

\begin{figure}[H]
	\centering
	\begin{tikzpicture}
		\begin{axis}[
				width=0.9\textwidth,
				height=0.42\textwidth,
				xlabel={Шаг валидации},
				ylabel={COMET},
				xmin=1,xmax=100,
				grid=both,
				major grid style={gray!30},
				minor grid style={gray!15},
			]
			\addplot[black, thick] table[x=step,y=comet,col sep=comma] {csv/final_validation.csv};
		\end{axis}
	\end{tikzpicture}
	\caption{Динамика COMET итоговой модели на validation-наборе}
	\label{fig:final-validation-comet}
\end{figure}

В начале обучения BLEU, chrF и COMET растут быстрее всего; затем графики переходят к более плавному участку. К концу обучения метрики в основном колеблются около достигнутого уровня. Поэтому простое добавление новых эпох без изменения данных, регуляризации или размера модели, скорее всего, дало бы меньший эффект, чем изменения на более ранних этапах экспериментов.

Итоговая оценка выполнялась на тестовом наборе, который не использовался при подборе гиперпараметров. На test-наборе модель получила следующие значения: \textbf{BLEU} -- 32{,}25, \textbf{chrF} -- 60{,}34, \textbf{COMET} -- 0{,}8519.

Примеры переводов итоговой модели приведены в таблице~\ref{tab:final-examples}.
В этих примерах модель корректно передает даты, в том числе века, записанные римскими цифрами, и сохраняет основные имена собственные.
Пунктуация и числовые фрагменты также в основном перенесены без искажений.

\begin{longtable}{|p{0.285\textwidth}|p{0.285\textwidth}|p{0.285\textwidth}|}
	\caption{Примеры переводов итоговой модели}\label{tab:final-examples}                                                                                                                                                                                                                                                                                                                                                                                                                                                                                                                                                           \\
	\hline
	\textbf{Исходное предложение}                                                                                                                                                                                     & \textbf{Образец}                                                                                                                                                                                      & \textbf{Перевод модели}                                                                                                                                                                             \\
	\hline
	Образец	\endfirsthead
	\hline
	\textbf{Исходное предложение}                                                                                                                                                                                     & \textbf{Образец}                                                                                                                                                                                      & \textbf{Перевод модели}                                                                                                                                                                             \\
	\hline
	\endhead
	Эпоху, в которой происходили данные события, принято называть Высоким Средневековьем, периодом европейской истории, затрагивающим XI, XII и XIII века (1000--1300 гг. н. э.).                                     & The age where the events took place is commonly referred as the High Middle Ages the period of European history in the 11th, 12th, and 13th centuries (AD 1000--1300).                                & The era in which these events occurred is called the High Middle Ages, a period of European history covering the 11th, 12th, and 13th centuries (1000--1300 AD).                                    \\
	\hline
	В 1683 году силы династии Цин (1644--1912) взяли под контроль западные и северные прибрежные районы Тайваня и в 1885 году объявили Тайвань провинцией империи Цин.                                                & In 1683, Qing dynasty (1644--1912) forces take control of Taiwan's western and northern coastal areas and declared Taiwan as a province of the Qing Empire in 1885.                                   & In 1683, the Qing Dynasty (1644--1912) took control of Taiwan's western and northern coastal areas and declared Taiwan a province of the Qing Empire in 1885.                                       \\
	\hline
	В 1994 году этот конфликт привёл к созданию на востоке Молдовы самопровозглашённой Приднестровской республики, которая имеет собственное правительство и валюту, но не признана ни одним государством-членом ООН. & In 1994, this conflict led to the creation of the self-proclaimed Transnistria Republic in eastern Moldova, which has its own government and currency but is not recognised by any UN member country. & In 1994, this conflict led to the creation of a self-proclaimed Transnistrian republic in eastern Moldova, which has its own government and currency, but is not recognized by any UN member state. \\
	\hline
\end{longtable}
